\documentclass[conference]{IEEEtran}
\usepackage[utf8]{inputenc}
\usepackage[T1]{fontenc}
\usepackage{cite}
\usepackage{amsmath,amssymb}
\usepackage{graphicx}
\usepackage{booktabs}
\usepackage{url}

\title{Evaluating a Generate--Validate--Repair Pipeline for LLM‑based Network Configuration: a Paired, Emulation‑Grounded Study}

\author{
\IEEEauthorblockN{Autonomous AI Researcher}
\IEEEauthorblockA{CNSM 2026 Agentic AI Researcher Track}
}

\begin{document}

\maketitle

\begin{abstract}
We preregistered and executed a paired, emulation‑grounded experiment (study id study-cand-1-gvr-verifier-adversarial-2026-08-01) to test whether integrating a deterministic network verifier into a generate--validate--repair (GVR) pipeline reduces the rate of unsafe configuration applications from an LLM relative to an LLM‑only pipeline. The preregistered primary estimand is the absolute paired difference in unsafe\_apply rate (guarded - baseline) where unsafe\_apply = 1 indicates a final configuration that would cause a semantic safety violation when applied and safe\_apply = 0 indicates no such violation. Execution used declared model gpt-5-mini and deterministic\_netops\_validator\_v5 under shared\_initial\_candidate semantics (exactly one initial LLM generation per task, reused across arms); guarded episodes could request up to one validator‑triggered repair. We ran N = 240 paired intent\ensuremath{\rightarrow}configuration tasks (144 in‑distribution, 48 distribution‑shift, 48 adversarial). All archived scoring artifacts report 'success' as a safe outcome; mapping: baseline\_success\_count = 240 corresponds to baseline unsafe\_apply = 0. Execution completed as planned (planned episodes = 480, completed = 480, hosted\_model\_call\_event\_count = 240, model\_calls\_used = 240). Every paired episode was scored safe under both arms (baseline = 240/240 safe; guarded = 240/240 safe). The paired risk difference (guarded - baseline) = 0.00; 95\% bootstrap percentile CI = [0.00, 0.00] (10,000 resamples, seed = 1729); exact McNemar two‑sided p = 1.0. Because shared\_initial\_candidate semantics produced identical shared candidates that passed deterministic validation in every task, no validator‑triggered repairs occurred and the paired contrast is degenerate. We report the preregistered design, artifact‑grounded execution accounting, full analysis mapping to archived artifacts, and an evidence‑based failure analysis that explains the null paired outcome and gives concrete recommendations for future informative GVR evaluations.
\end{abstract}

\section{Introduction}
Automating the translation of high‑level operator intent into device configurations is a central NetOps goal. Large language models (LLMs) and agentic workflows increasingly propose configuration change plans, but token‑level correctness metrics do not capture ordering, dependency, and invariants needed for safe application in live networks. Closed‑loop, emulation‑grounded evaluation---where generated configs are applied to device images or simulators and deterministically scored---reveals semantic failure modes (ordering mistakes, transient violations, protected‑resource touches) not visible in text‑only tests [1].

A common architectural remedy is generate--validate--repair (GVR): an LLM generates a candidate configuration, a deterministic verifier checks semantic invariants, and an automated repair routine rewrites the candidate when violations are detected [2,3]. Prior work shows GVR‑style patterns are promising for syntactic and many semantic faults, but open questions remain about empirical verifier benefit under distribution shift and adversarial prompts, and about experimental designs that reliably expose paired differences between guarded and unguarded pipelines [2,4].

This paper reports a fully preregistered paired experiment that measured a single, precisely defined estimand: the absolute paired difference in unsafe\_apply rate (guarded minus baseline), executed under shared\_initial\_candidate semantics (one initial shared LLM generation per task reused across both arms). We preserve and reference the archived execution and analysis artifacts throughout the manuscript. Our reported contributions are: (1) an artifact‑grounded execution and analysis of a preregistered paired GVR test (study id study-cand-1-gvr-verifier-adversarial-2026-08-01); (2) explicit mapping between the preregistered estimand and archived analysis outputs (contingency CSVs, bootstrap parameters, and reconciliation artifacts); and (3) a failure analysis explaining the null paired outcome and concrete artifact‑grounded recommendations for future GVR evaluations.

\section{Related Work}
Executable NetOps benchmarks and emulation frameworks have recently been proposed to surface workflow and planning faults for LLM agents rather than only assessing textual fidelity; NetConfArena is an example of this executable, closed‑loop approach [1]. Work on automatic configuration repair and program repair documents generate--validate patterns and demonstrates that validators often enable high repair coverage for many benchmarked faults, although semantic coverage varies by verifier design [2,4]. Survey and position papers synthesize agentic architectures, tool augmentation, verifier integration, and uncertainty calibration as promising patterns for operational NetOps automation and safety [3,5]. Evolution of configuration analysis tools such as Batfish provides engineering precedent for verifier design and incremental verification [6].

Our study differs from prior evaluations by combining: (a) a predeclared paired estimand; (b) preregistered shared\_initial\_candidate execution semantics; (c) emulation‑grounded deterministic scoring; and (d) full archival of per‑call traces, validator traces, and analysis artifacts so that the exact measured estimand can be reconciled to archived outputs.

\section{Methodology}
Preregistration, estimand, and mapping: The experiment was preregistered (study id study-cand-1-gvr-verifier-adversarial-2026-08-01; preregistration SHA‑256 recorded in the bundle). Primary estimand: paired\_success\_rate\_difference\_guarded\_minus\_baseline defined as the absolute paired difference in unsafe\_apply rate (guarded - baseline) where unsafe\_apply = 1 indicates a final configuration that would cause a semantic safety violation when applied; safe\_apply = 0 indicates no such violation. The archived analysis executor reports 'success' as the safe outcome in the archived condition summary; we explicitly map archived fields to the estimand: archived baseline\_success\_count = 240 equals baseline\_safe\_apply\_count = 240, hence baseline\_unsafe\_apply\_count = 0. The mapping between preregistered estimand names and the analysis artifact names is summarized here and archived: primary estimand \ensuremath{\rightarrow} the archived paired contingency table and the archived condition summary; bootstrap parameters \ensuremath{\rightarrow} the archived analysis artifact (bootstrap\_resamples = 10000, bootstrap\_seed = 1729). This mapping ensures that the numeric outputs reported here are traceable to archived artifacts.

Design, tasks, and sampling: The design was paired with shared\_initial\_candidate generation semantics (the adapter produced exactly one initial LLM generation per task and both arms used that same generation). N = 240 distinct intent\ensuremath{\rightarrow}configuration tasks were blocked into topology/vendor clusters and composed as: in‑distribution 60\% (144 tasks), distribution‑shift 20\% (48), and adversarial‑prompt 20\% (48). Tasks were generated deterministically by deterministic\_netops\_task\_generator\_v5 and include intent payloads, initial states, required command sequences, and preserved‑field constraints (representative task manifests are archived in the archived task manifest).

Model, oracle, and adapter contract: Execution used the declared model gpt-5-mini (the archived model configuration; requested\_model = "gpt-5-mini"). The deterministic verifier was deterministic\_netops\_validator\_v5; the guarded pipeline could request up to one validator‑triggered repair per task. The adapter family was hosted\_netops\_gvr\_v1 and execution used the registered shared\_initial\_candidate semantics: exactly one shared initial generation was produced per task (provider\_call\_audit.hosted\_model\_call\_event\_count = 240). The preregistered adapter contract constrains interpretation: under shared\_initial\_candidate semantics, paired differences can only arise when the single shared candidate fails validation or when the guarded arm applies a repair and thereby changes the final configuration relative to the baseline.

Scoring, exclusions, and missingness: The deterministic scorer assigns a binary final outcome: safe (0) or unsafe (1) based on emulator execution and semantic invariant checks recorded in the validator trace (example: the archived execution artifact). Preregistered exclusions: pairs flagged as 'possible false-safe' (validator PASS but emulator unsafe outcome) are excluded from primary analysis and handled in sensitivity analyses; infra failures after allowed retries were to be marked infra‑failed and excluded per preregistration. The retry policy permitted up to 2 infra‑retries; no infra failures occurred in this run. All missingness and exclusion counts are archived (the archived analysis artifact).

Statistical analysis: The primary analysis executor was paired\_binary\_analysis\_v1 and produced a paired contingency table (the archived paired contingency table), bootstrap 95\% percentile CI using 10,000 resamples (bootstrap\_seed = 1729), and an exact two‑sided McNemar test on discordant pairs. Secondary planned outputs included median repair counts and model‑call budgets per task; repair\_attempts = 0 across archived repair logs. Multiplicity: single primary hypothesis; secondary tests would be FDR‑corrected when reported. All analysis parameter values and logs are archived (the archived analysis artifact).

\section{Results}
Execution accounting and audit: The run completed per the preregistered plan. Planned\_episode\_count = 480; completed\_episode\_count = 480; complete\_pair\_count = 240. Provider call audit records hosted\_model\_call\_event\_count = 240 and cache\_miss\_count = 240 consistent with exactly one shared\_initial\_generation per task (deterministic\_reconciliation.json). Model\_calls\_used = 240 (execution\_manifest). No infra failures, no unscorable episodes, and no preregistered exclusions occurred (the archived analysis artifact).

Primary binary outcomes and contingency: Every paired episode was scored safe under both arms. the archived condition summary records baseline successes = 240, baseline failures = 0 (baseline unsafe\_apply = 0) and guarded successes = 240, guarded failures = 0 (guarded unsafe\_apply = 0). The paired contingency (the archived paired contingency table) is n11 (both safe) = 240; n10 = 0; n01 = 0; n00 = 0. Hence the paired estimator (guarded\_unsafe\_rate - baseline\_unsafe\_rate) = 0.00. Bootstrap percentile 95\% CI = [0.00, 0.00] (resamples = 10,000; seed = 1729) and exact McNemar two‑sided p = 1.0. Analysis artifacts (the archived analysis results, CSVs, and analysis\_log.jsonl) are archived and referenced by the execution manifest.

Validator activity and repair budget utilization: No validator‑triggered repairs were applied (repair\_attempts = 0 across the repair log). Representative validator traces (e.g., the archived execution artifact) show validation\_before and validation\_after fields, final\_state snapshots, required\_command\_count, difficulty labels ("very\_hard"/"extreme"), and validator PASS semantics. Because the guarded final configuration matched the shared initial candidate for every episode, baseline and guarded outcomes are identical across all tasks and subgroups (in‑distribution, shift, adversarial all show zero unsafe applies). The reconciliation artifact deterministic\_reconciliation.json records all\_deterministic\_consistency\_checks\_passed = true and provider\_call\_audit.stage\_counts.shared\_initial\_generation = 240, confirming contract adherence.

Planned but unused sensitivity channels: The preregistration specified sensitivity analyses for 'possible false‑safe' flags and infra failures (treating those items as unsafe in worst‑case analyses). No items were flagged; contamination\_summary.csv is empty. The planned subgroup analyses (per‑condition paired contrasts for in‑distribution, shift, adversarial) are degenerate here (all zeros) and the archived subgroup contingency slices are available in the archived paired contingency table for independent inspection.

\section{Discussion}
Why the paired contrast is null (artifact‑grounded explanation): The preregistered shared\_initial\_candidate semantics produced exactly one shared LLM output per task (provider\_call\_audit), and deterministic\_netops\_validator\_v5 passed every shared candidate (the archived execution artifact). Under the adapter contract the guarded arm could only change the paired outcome by invoking a validator‑triggered repair; because no repair was invoked the guarded arm could not create discordant pairs. This mechanistic, artifact‑grounded explanation accounts for the degenerate paired contingency and is fully consistent with the archived analysis executor output (estimate = 0.0, CI = [0.0, 0.0], p = 1.0).

Design tradeoffs and practical implications: The shared\_initial\_candidate semantics reduces model‑call budget and ensures a directly paired LLM sample, but it also prevents generator variability between arms and therefore makes empirical verifier benefit contingent on initial candidate failures. To measure verifier effectiveness as a causal intervention one must either (A) allow independent initial generations per arm, (B) introduce controlled generator perturbations or higher adversary strength so validators have opportunities to flag failures, or (C) expand repair budgets so the repair mechanism can be exercised when validation fails. These recommendations are directly motivated by provider\_call\_audit and validator traces showing identical shared\_initial responses and zero repair usage.

Threats to validity (expanded and artifact‑grounded): Internal validity: the shared\_initial\_candidate contract is an explicit design choice that restricts causal contrast interpretation; all evidence of that contract and its runtime adherence is archived (the archived model configuration and deterministic\_reconciliation.json). Construct validity: deterministic\_netops\_validator\_v5 encodes a finite check set; while validator traces show comprehensive recording of observed\_touched\_field\_count and violation\_codes, uncovered semantic fault classes may remain. External validity: tasks used a NetConfArena‑style synthetic emulation corpus and deterministic task generator (deterministic\_netops\_task\_generator\_v5); extrapolation to production multi‑vendor deployments requires cautious qualification. Conclusion validity: realized baseline unsafe rate = 0 prevents detection of verifier benefit in this design; future power considerations must account for generator variability and adversary strength to ensure informative discordant pairs.

Recommendations for future paired GVR evaluations (artifact‑grounded): (1) If paired causal estimation of verifier benefit is the goal, avoid shared\_initial\_candidate semantics or intentionally inject generator variability for one arm (provider\_call\_audit demonstrates effect of single vs multiple generations). (2) Increase repair budget (allow multiple repair iterations) and instrument repair prompts/edits in repair\_provider\_trace to study repair search behavior. (3) Expand adversary protocols beyond single fixed seed adversaries to probe adaptive prompt manipulations; archived contamination and adversarial tagging logic provide a starting point for reproducible adversary generation. (4) Improve verifier semantic coverage and add cross‑checks between validator PASS and emulator outcomes to detect potential false‑safe cases; the preregistered contamination plan included these checks and their automation can be extended.

\section{Limitations}
\begin{itemize}
\item Emulator and corpus scope: tasks used a NetConfArena‑style synthetic emulation corpus; external validity to production, multi‑vendor networks and operator workflows is limited.
\item Single model: a single hosted model (gpt‑5‑mini) was tested; results do not imply generality across models, prompts, or prompting strategies.
\item Shared\_initial\_candidate semantics: the paired design used one shared initial generation per task; this prevents generator variability between arms and makes verifier benefit contingent on initial candidate failures.
\item Adversary bounds: adversarial prompts were generated under a preregistered, fixed adversary protocol and do not represent unconstrained, adaptive adversaries.
\item Verifier coverage: deterministic\_netops\_validator\_v5 implements a finite set of checks; uncovered semantic fault classes may remain undetected by the verifier, and verifier PASS vs simulator outcome mismatches were planned to be flagged and analyzed (none occurred in this run).
\end{itemize}

\section{Conclusion}
We executed a fully preregistered, paired, emulation‑grounded evaluation of a GVR pipeline for LLM‑based network configuration. Observed outcome: every paired episode was safe under both arms, yielding a paired difference = 0.00 (95\% bootstrap CI [0.00, 0.00]; McNemar p = 1.0). Importantly---and stated explicitly per reviewer request---under the preregistered shared\_initial\_candidate semantics the paired contrast could only differ if the single shared initial candidate failed deterministic validation or if a validator‑triggered repair was applied; in this run the shared initial candidates passed validation and no repairs were invoked, producing the degenerate paired outcome. Our artifact archive (preregistration, per‑call traces, validator traces, and analysis artifacts) supports these conclusions and provides concrete guidance for future experiments that seek informative paired contrasts by varying generation semantics, adversary strength, or repair budgets.

References
[1] C. Liu, X. Xie, X. Chen, and Y. Cui, "NetConfArena: An Executable Benchmark for LLM Agents in Closed‑Loop Network Configuration," arXiv:2608.23179, 2026.
[2] Z. Gu, P. Zhang, X. Feng, and X. Liu, "Astragalus: Automatic Configuration Repair for Production Networks," arXiv:2605.22092, 2026.
[3] M. Bilal, J. Crowcroft, R. Wang, X. Xu, and S. Dustdar, "Large Language Models for Agentic NetOps and AIOps: Architectures, Evaluation, and Safety," arXiv:2605.12729, 2026.
[4] A. Yehudai et al., "A Survey on Evaluation of LLM‑based Agents," Findings of ACL, 2026, 10.18653/v1/2026.findings-acl.1330.
[5] N. Van Tu, S. Nam, and J. W.‑K. Hong, "Intent‑Based Network Configuration Using Large Language Models," Netw. Model. Anal. Health Inf. Bioinform., 2024, 10.1002/nem.2313.
[6] M. Brown et al., "Lessons from the evolution of the Batfish configuration analysis tool," Proc. ACM (conference record), 2023, 10.1145/3603269.3604866.
[7] W. Shi et al., "EHRAgent: Code Empowers Large Language Models for Few‑shot Complex Tabular Reasoning on Electronic Health Records," in Proc. EMNLP, 2023, 10.18653/v1/2023.emnlp-main.1245.

Figure captions
Figure 1 --- Paired unsafe‑apply rates by subgroup (in‑distribution, distribution‑shift, adversarial): all subgroups show 0\% unsafe applies for both baseline and guarded arms (the archived condition summary).

Figure 2 --- Repair attempts per guarded episode (boxplot): zero repair invocations observed across all guarded episodes (repair\_attempts = 0 in archived repair logs).

Table captions
Table 1 --- Paired contingency table (guarded vs baseline): n11 = 240, n10 = 0, n01 = 0, n00 = 0; paired difference = 0.00 (95\% bootstrap CI [0.00,0.00], bootstrap\_resamples = 10000, seed = 1729).

Table 2 --- Task composition and counts: in‑distribution 144, distribution‑shift 48, adversarial 48; total paired tasks 240; planned episodes 480, completed 480, hosted\_model\_call\_event\_count = 240.

Disclosure Statement
Models and orchestration: Execution used a single hosted model (declared\_model\_version = "gpt-5-mini") via the registered adapter hosted\_netops\_gvr\_v1. The guarded pipeline used deterministic\_netops\_validator\_v5 and permitted up to one validator‑triggered repair per task. Runtime sampling: the archived model configuration records temperature = null (unset); unset temperature does not imply determinism by the hosted provider and sampling nondeterminism remains possible; all model‑call metadata were archived. Preregistration and autonomy: the study was preregistered (study id study-cand-1-gvr-verifier-adversarial-2026-08-01; preregistration SHA‑256: ) and executed under the frozen plan after the autonomy lock. Human role and autonomy boundary: prior to the autonomy lock the human Research Director selected the topic family and constrained the initial master prompt; all literature search, experiment selection, preregistration, execution, analysis, and manuscript drafting occurred autonomously after the lock per the track rules. Preregistered generation semantics: the registered adapter used shared\_initial\_candidate semantics (one shared initial generation per task reused across both arms); therefore only validator‑triggered repairs (or validator failures of the shared candidate) could create paired divergences; this constraint is reported throughout Methods, Results, Abstract, and Conclusion. Execution accounting (archived): planned\_episode\_count = 480, completed\_episode\_count = 480, complete\_pair\_count = 240, hosted\_model\_call\_event\_count = 240, model\_calls\_used = 240. Immutable master prompt: the initial master prompt is archived at the archived provenance artifact (SHA‑256: ). Reproducibility pointer: archived execution manifest, per‑call traces, task manifests, validator traces, and analysis artifacts are included in the execution bundle and referenced in the manuscript (the archived analysis artifacts.csv, the archived execution artifacts.jsonl). No manual human editing of the technical manuscript content occurred after the autonomy lock as required by the track rules.

(End of manuscript.)

\section*{Disclosure Statement}
Execution used declared model gpt-5-mini via the registered adapter hosted\_netops\_gvr\_v1; guarded pipeline used deterministic\_netops\_validator\_v5 and permitted <=1 validator‑triggered repair per task. The study was preregistered (study id study-cand-1-gvr-verifier-adversarial-2026-08-01; preregistration SHA‑256: ) and executed under the frozen plan after the autonomy lock. The human Research Director selected the topic family and constrained the initial master prompt prior to the lock; all literature search, experiment design, preregistration, execution, analysis, and manuscript drafting occurred autonomously after the lock. The registered adapter used shared\_initial\_candidate semantics (one shared initial generation per task reused across both arms); under that contract the paired contrast could only differ if the shared initial candidate failed deterministic validation or if a validator‑triggered repair was applied. Immutable master prompt archived at provenance/master\_prompt.txt (SHA‑256: 1872df1e\allowbreak{}1805d2d9\allowbreak{}6940456c\allowbreak{}a016bd66\allowbreak{}5d1d5196\allowbreak{}add77f5a\allowbreak{}cdf1582b\allowbreak{}b39b15ba). Execution artifacts (the archived execution artifact, the archived execution artifact, the archived task manifest, the archived analysis artifacts.csv, the archived analysis results) are archived and referenced in the manuscript. No manual human editing of technical manuscript content occurred post‑lock.

\begin{thebibliography}{99}
\bibitem{ref1} https://arxiv.org/abs/2608.23179
\bibitem{ref2} https://arxiv.org/abs/2605.22092
\bibitem{ref3} https://arxiv.org/abs/2605.12729
\bibitem{ref4} https://doi.org/10.18653/v1/2026.findings-acl.1330
\bibitem{ref5} https://doi.org/10.1002/nem.2313
\bibitem{ref6} https://doi.org/10.1145/3603269.3604866
\bibitem{ref7} https://doi.org/10.18653/v1/2023.emnlp-main.1245
\end{thebibliography}

\end{document}
